% Options for packages loaded elsewhere
\PassOptionsToPackage{unicode}{hyperref}
\PassOptionsToPackage{hyphens}{url}
\documentclass[
]{article}
\usepackage{xcolor}
\usepackage[margin=1in]{geometry}
\usepackage{amsmath,amssymb}
\setcounter{secnumdepth}{-\maxdimen} % remove section numbering
\usepackage{iftex}
\ifPDFTeX
  \usepackage[T1]{fontenc}
  \usepackage[utf8]{inputenc}
  \usepackage{textcomp} % provide euro and other symbols
\else % if luatex or xetex
  \usepackage{unicode-math} % this also loads fontspec
  \defaultfontfeatures{Scale=MatchLowercase}
  \defaultfontfeatures[\rmfamily]{Ligatures=TeX,Scale=1}
\fi
\usepackage{lmodern}
\ifPDFTeX\else
  % xetex/luatex font selection
\fi
% Use upquote if available, for straight quotes in verbatim environments
\IfFileExists{upquote.sty}{\usepackage{upquote}}{}
\IfFileExists{microtype.sty}{% use microtype if available
  \usepackage[]{microtype}
  \UseMicrotypeSet[protrusion]{basicmath} % disable protrusion for tt fonts
}{}
\makeatletter
\@ifundefined{KOMAClassName}{% if non-KOMA class
  \IfFileExists{parskip.sty}{%
    \usepackage{parskip}
  }{% else
    \setlength{\parindent}{0pt}
    \setlength{\parskip}{6pt plus 2pt minus 1pt}}
}{% if KOMA class
  \KOMAoptions{parskip=half}}
\makeatother
\usepackage{color}
\usepackage{fancyvrb}
\newcommand{\VerbBar}{|}
\newcommand{\VERB}{\Verb[commandchars=\\\{\}]}
\DefineVerbatimEnvironment{Highlighting}{Verbatim}{commandchars=\\\{\}}
% Add ',fontsize=\small' for more characters per line
\usepackage{framed}
\definecolor{shadecolor}{RGB}{248,248,248}
\newenvironment{Shaded}{\begin{snugshade}}{\end{snugshade}}
\newcommand{\AlertTok}[1]{\textcolor[rgb]{0.94,0.16,0.16}{#1}}
\newcommand{\AnnotationTok}[1]{\textcolor[rgb]{0.56,0.35,0.01}{\textbf{\textit{#1}}}}
\newcommand{\AttributeTok}[1]{\textcolor[rgb]{0.13,0.29,0.53}{#1}}
\newcommand{\BaseNTok}[1]{\textcolor[rgb]{0.00,0.00,0.81}{#1}}
\newcommand{\BuiltInTok}[1]{#1}
\newcommand{\CharTok}[1]{\textcolor[rgb]{0.31,0.60,0.02}{#1}}
\newcommand{\CommentTok}[1]{\textcolor[rgb]{0.56,0.35,0.01}{\textit{#1}}}
\newcommand{\CommentVarTok}[1]{\textcolor[rgb]{0.56,0.35,0.01}{\textbf{\textit{#1}}}}
\newcommand{\ConstantTok}[1]{\textcolor[rgb]{0.56,0.35,0.01}{#1}}
\newcommand{\ControlFlowTok}[1]{\textcolor[rgb]{0.13,0.29,0.53}{\textbf{#1}}}
\newcommand{\DataTypeTok}[1]{\textcolor[rgb]{0.13,0.29,0.53}{#1}}
\newcommand{\DecValTok}[1]{\textcolor[rgb]{0.00,0.00,0.81}{#1}}
\newcommand{\DocumentationTok}[1]{\textcolor[rgb]{0.56,0.35,0.01}{\textbf{\textit{#1}}}}
\newcommand{\ErrorTok}[1]{\textcolor[rgb]{0.64,0.00,0.00}{\textbf{#1}}}
\newcommand{\ExtensionTok}[1]{#1}
\newcommand{\FloatTok}[1]{\textcolor[rgb]{0.00,0.00,0.81}{#1}}
\newcommand{\FunctionTok}[1]{\textcolor[rgb]{0.13,0.29,0.53}{\textbf{#1}}}
\newcommand{\ImportTok}[1]{#1}
\newcommand{\InformationTok}[1]{\textcolor[rgb]{0.56,0.35,0.01}{\textbf{\textit{#1}}}}
\newcommand{\KeywordTok}[1]{\textcolor[rgb]{0.13,0.29,0.53}{\textbf{#1}}}
\newcommand{\NormalTok}[1]{#1}
\newcommand{\OperatorTok}[1]{\textcolor[rgb]{0.81,0.36,0.00}{\textbf{#1}}}
\newcommand{\OtherTok}[1]{\textcolor[rgb]{0.56,0.35,0.01}{#1}}
\newcommand{\PreprocessorTok}[1]{\textcolor[rgb]{0.56,0.35,0.01}{\textit{#1}}}
\newcommand{\RegionMarkerTok}[1]{#1}
\newcommand{\SpecialCharTok}[1]{\textcolor[rgb]{0.81,0.36,0.00}{\textbf{#1}}}
\newcommand{\SpecialStringTok}[1]{\textcolor[rgb]{0.31,0.60,0.02}{#1}}
\newcommand{\StringTok}[1]{\textcolor[rgb]{0.31,0.60,0.02}{#1}}
\newcommand{\VariableTok}[1]{\textcolor[rgb]{0.00,0.00,0.00}{#1}}
\newcommand{\VerbatimStringTok}[1]{\textcolor[rgb]{0.31,0.60,0.02}{#1}}
\newcommand{\WarningTok}[1]{\textcolor[rgb]{0.56,0.35,0.01}{\textbf{\textit{#1}}}}
\usepackage{graphicx}
\makeatletter
\newsavebox\pandoc@box
\newcommand*\pandocbounded[1]{% scales image to fit in text height/width
  \sbox\pandoc@box{#1}%
  \Gscale@div\@tempa{\textheight}{\dimexpr\ht\pandoc@box+\dp\pandoc@box\relax}%
  \Gscale@div\@tempb{\linewidth}{\wd\pandoc@box}%
  \ifdim\@tempb\p@<\@tempa\p@\let\@tempa\@tempb\fi% select the smaller of both
  \ifdim\@tempa\p@<\p@\scalebox{\@tempa}{\usebox\pandoc@box}%
  \else\usebox{\pandoc@box}%
  \fi%
}
% Set default figure placement to htbp
\def\fps@figure{htbp}
\makeatother
\setlength{\emergencystretch}{3em} % prevent overfull lines
\providecommand{\tightlist}{%
  \setlength{\itemsep}{0pt}\setlength{\parskip}{0pt}}
\usepackage{bookmark}
\IfFileExists{xurl.sty}{\usepackage{xurl}}{} % add URL line breaks if available
\urlstyle{same}
\hypersetup{
  pdftitle={Midterm 1 W26},
  pdfauthor={Darius Arasteh},
  hidelinks,
  pdfcreator={LaTeX via pandoc}}

\title{Midterm 1 W26}
\author{Darius Arasteh}
\date{2026-01-29}

\begin{document}
\maketitle

\subsection{Instructions}\label{instructions}

Answer the following questions and complete the exercises in RMarkdown.
Please embed all of your code and push your final work to your
repository. Your code must be organized, clean, and run free from
errors. Remember, you must remove the \texttt{\#} for any included code
chunks to run. Be sure to add your name to the author header above.

Your code must knit in order to be considered. If you are stuck and
cannot answer a question, then comment out your code and knit the
document. You may use your notes, labs, and homework to help you
complete this exam. Do not use any other resources- including AI
assistance or other students' work.

Don't forget to answer any questions that are asked in the prompt! Each
question must be coded; it cannot be answered by a sort in a spreadsheet
or a written response only.

For all plots you create, a title and clearly labeled axes must be
provided. We also expect pipes \texttt{\%\textgreater{}\%} to be used
wherever possible.

Be sure to push your completed midterm to your repository and upload the
document to Gradescope. This exam is worth 50 points.

Please load the following libraries.

\begin{Shaded}
\begin{Highlighting}[]
\FunctionTok{library}\NormalTok{(tidyverse)}
\FunctionTok{library}\NormalTok{(janitor)}
\end{Highlighting}
\end{Shaded}

\subsection{Part 1: Repository}\label{part-1-repository}

\textbf{Question 1. (3 points) Before you start analyzing data, please
put a link to your GitHub repository below. Your repository should have
a clear README and be well-organized. Add \texttt{jmledford3115} and
\texttt{bryshalm} as collaborators to your repository if you haven't
already done so.}

Link to repository:\\
\url{https://github.com/darasteh-gif/BIS15L_W26_darasteh}\strut \\
\strut \\
\strut \\
\strut \\

\subsection{Part 2: Data and Analysis}\label{part-2-data-and-analysis}

In the midterm 1 folder there is a second folder called \texttt{data}.
Inside the \texttt{data} folder, there is a .csv file called
\texttt{anolis\_dat.csv}. These data came from
\href{https://academic.oup.com/evolut/article/64/9/2731/6854302?login=true}{D.
Luke Mahler, Liam J. Revell, Richard E. Glor, Jonathan B. Losos,
ECOLOGICAL OPPORTUNITY AND THE RATE OF MORPHOLOGICAL EVOLUTION IN THE
DIVERSIFICATION OF GREATER ANTILLEAN ANOLES, Evolution, Volume 64, Issue
9, 1 September 2010, Pages 2731--2745}. The original research article is
included in the data folder.

\emph{Anolis} is a genus of lizards commonly known as anoles. Anoles are
found throughout the Americas, but are especially diverse in the
Caribbean. The data include morphological measurements for \emph{Anolis}
lizards from the islands of the Greater Antilles. These data can be used
to study patterns of morphological evolution and adaptation in
\emph{Anolis} lizards.

The variables include:\\
- \texttt{species}: Species name of the anole lizard.\\
- \texttt{habitat}: Habitat type where the lizard was found.\\
- \texttt{hindlimb\_length\_mm}: Length of the lizard's hindlimbs (in
millimeters).\\
- \texttt{tail\_length\_mm}: Length of the lizard's tail (in
millimeters).\\
- \texttt{body\_length\_mm}: Length of the lizard's body (in
millimeters).\\
- \texttt{toepad\_lamellae\_count}: Count of lamellae on the lizard's
toepads.\\
- \texttt{island}: Island where the lizard was found.

\textbf{Question 2. (2 points) Load the data and store it as an object
called \texttt{anolis}.}\\
\strut \\

\begin{Shaded}
\begin{Highlighting}[]
\NormalTok{anolis }\OtherTok{\textless{}{-}} \FunctionTok{read.csv}\NormalTok{(}\StringTok{"data/anolis\_dat.csv"}\NormalTok{) }\SpecialCharTok{\%\textgreater{}\%} \FunctionTok{clean\_names}\NormalTok{()}
\end{Highlighting}
\end{Shaded}

\hfill\break
\hfill\break
\hfill\break

\textbf{Question 3. (2 points) Use a summary function of your choice to
get an idea of the structure of the data.}\\

\begin{Shaded}
\begin{Highlighting}[]
\FunctionTok{summary}\NormalTok{(anolis)}
\end{Highlighting}
\end{Shaded}

\begin{verbatim}
##    species            habitat          hindlimb_length_mm tail_length_mm  
##  Length:52          Length:52          Min.   :20.00      Min.   : 39.79  
##  Class :character   Class :character   1st Qu.:31.70      1st Qu.: 82.19  
##  Mode  :character   Mode  :character   Median :40.01      Median :101.86  
##                                        Mean   :40.51      Mean   :102.81  
##                                        3rd Qu.:49.27      3rd Qu.:121.92  
##                                        Max.   :65.48      Max.   :154.45  
##  body_length_mm  toepad_lamellae_count    island         
##  Min.   :29.01   Min.   :23.00         Length:52         
##  1st Qu.:40.07   1st Qu.:28.00         Class :character  
##  Median :47.05   Median :29.00         Mode  :character  
##  Mean   :47.65   Mean   :30.54                           
##  3rd Qu.:54.53   3rd Qu.:33.00                           
##  Max.   :72.32   Max.   :44.00
\end{verbatim}

\hfill\break
\hfill\break
\hfill\break
\hfill\break

\textbf{Question 4. (2 points) Clean the variable names so they are all
lowercase and without special characters or spaces. Be sure to use the
cleaned data for all subsequent analyses.}\\
\textbf{I Cleaned the names when I loaded the data}

\hfill\break
\hfill\break
\hfill\break

\textbf{Question 5. (4 points) Convert the \texttt{habitat} and
\texttt{island} variables to factors.}\\
\strut \\

\begin{Shaded}
\begin{Highlighting}[]
\NormalTok{anolis }\SpecialCharTok{\%\textgreater{}\%} 
  \FunctionTok{mutate}\NormalTok{(}\AttributeTok{habitat=}\FunctionTok{as.factor}\NormalTok{(habitat)) }\SpecialCharTok{\%\textgreater{}\%} 
  \FunctionTok{mutate}\NormalTok{(}\AttributeTok{habitat=}\FunctionTok{as.factor}\NormalTok{(habitat)) }
\end{Highlighting}
\end{Shaded}

\begin{verbatim}
##               species      habitat hindlimb_length_mm tail_length_mm
## 1             A. ahli Trunk-ground              50.46          81.99
## 2          A. alayoni         Twig              25.50          54.75
## 3          A. alfaroi   Grass-bush              26.17          79.00
## 4         A. aliniger  Trunk-crown              36.80          84.88
## 5         A. allisoni  Trunk-crown              50.39         154.45
## 6          A. allogus Trunk-ground              49.37          91.01
## 7          A. alumina   Grass-bush              29.97         106.40
## 8        A. alutaceus   Grass-bush              27.39          94.62
## 9      A. angusticeps         Twig              24.37          65.10
## 10         A. armouri Trunk-ground              51.81         101.37
## 11   A. bahorucoensis   Grass-bush              38.71         115.82
## 12         A. bremeri Trunk-ground              47.44         114.01
## 13        A. breslini Trunk-ground              49.84         122.85
## 14    A. chlorocyanus  Trunk-crown              48.55         142.12
## 15       A. clivicola   Grass-bush              34.79          98.75
## 16        A. confusus Trunk-ground              41.63          78.31
## 17           A. cooki Trunk-ground              48.50         134.57
## 18    A. cupeyalensis   Grass-bush              22.88          93.90
## 19    A. cyanopleurus   Grass-bush              31.38         111.68
## 20 A. dolichocephalus   Grass-bush              39.49         131.27
## 21         A. grahami  Trunk-crown              48.71         119.78
## 22           A. guafe Trunk-ground              39.09          69.45
## 23         A. guazuma         Twig              20.46          39.79
## 24       A. gundlachi Trunk-ground              58.32         136.13
## 25       A. haetianus Trunk-ground              65.48         154.41
## 26      A. hendersoni   Grass-bush              38.53         128.81
## 27      A. homolechis Trunk-ground              44.97          90.92
## 28           A. imias Trunk-ground              51.67         104.67
## 29    A. inexpectatus   Grass-bush              26.80         101.12
## 30       A. insolitus         Twig              26.13          50.77
## 31           A. jubar Trunk-ground              41.08          82.15
## 32           A. krugi   Grass-bush              40.53         130.76
## 33      A. lineatopus Trunk-ground              52.71         119.97
## 34   A. longitibialis Trunk-ground              61.92         140.36
## 35      A. macilentus   Grass-bush              31.80          97.35
## 36        A. marcanoi Trunk-ground              49.24         123.95
## 37         A. mestrei Trunk-ground              44.93          92.11
## 38        A. opalinus  Trunk-crown              35.67          82.15
## 39      A. ophiolepis   Grass-bush              26.88          81.13
## 40        A. placidus         Twig              20.00          47.96
## 41       A. poncensis   Grass-bush              33.39         118.20
## 42        A. porcatus  Trunk-crown              44.58         140.26
## 43 A. quadriocellifer Trunk-ground              39.05          82.20
## 44     A. rubribarbus Trunk-ground              50.43          93.94
## 45          A. sagrei Trunk-ground              45.47         111.90
## 46    A. semilineatus   Grass-bush              33.99         121.61
## 47         A. shrevei Trunk-ground              43.61         102.53
## 48      A. singularis  Trunk-crown              38.42         102.34
## 49         A. strahmi Trunk-ground              65.24         138.87
## 50       A. stratulus  Trunk-crown              34.38          77.93
## 51       A. vanidicus   Grass-bush              26.45          90.60
## 52       A. whitemani Trunk-ground              51.06         115.26
##    body_length_mm toepad_lamellae_count      island
## 1           51.67                    27        Cuba
## 2           41.32                    31        Cuba
## 3           30.95                    24        Cuba
## 4           51.53                    36  Hispaniola
## 5           72.32                    41        Cuba
## 6           51.72                    28        Cuba
## 7           32.94                    29  Hispaniola
## 8           31.84                    28        Cuba
## 9           40.22                    28        Cuba
## 10          56.11                    31  Hispaniola
## 11          41.81                    32  Hispaniola
## 12          55.65                    32        Cuba
## 13          52.29                    29  Hispaniola
## 14          65.44                    42  Hispaniola
## 15          39.03                    26        Cuba
## 16          46.72                    32        Cuba
## 17          54.45                    33 Puerto Rico
## 18          29.01                    24        Cuba
## 19          34.32                    26        Cuba
## 20          45.34                    36  Hispaniola
## 21          57.97                    38     Jamacia
## 22          43.95                    29        Cuba
## 23          39.24                    28        Cuba
## 24          59.97                    28 Puerto Rico
## 25          68.19                    33  Hispaniola
## 26          43.19                    34  Hispaniola
## 27          51.34                    32        Cuba
## 28          54.89                    28        Cuba
## 29          31.28                    27        Cuba
## 30          40.70                    25  Hispaniola
## 31          47.38                    29        Cuba
## 32          44.35                    33 Puerto Rico
## 33          56.50                    29     Jamacia
## 34          63.29                    28  Hispaniola
## 35          37.39                    28        Cuba
## 36          53.79                    33  Hispaniola
## 37          49.05                    30        Cuba
## 38          42.27                    35     Jamacia
## 39          34.59                    24        Cuba
## 40          39.63                    26  Hispaniola
## 41          41.52                    29 Puerto Rico
## 42          64.37                    44        Cuba
## 43          45.03                    28        Cuba
## 44          53.74                    29        Cuba
## 45          53.14                    33        Cuba
## 46          36.68                    30  Hispaniola
## 47          48.85                    28  Hispaniola
## 48          52.65                    38  Hispaniola
## 49          65.36                    31  Hispaniola
## 50          43.62                    34 Puerto Rico
## 51          34.19                    23        Cuba
## 52          54.77                    29  Hispaniola
\end{verbatim}

\hfill\break
\hfill\break
\hfill\break

\textbf{Question 6. (2 points) Anole species were sampled from multiple
islands. Which islands are represented in the data? Display the island
names.}

\begin{Shaded}
\begin{Highlighting}[]
\NormalTok{anolis }\SpecialCharTok{\%\textgreater{}\%} 
  \FunctionTok{distinct}\NormalTok{(island)}
\end{Highlighting}
\end{Shaded}

\begin{verbatim}
##        island
## 1        Cuba
## 2  Hispaniola
## 3 Puerto Rico
## 4     Jamacia
\end{verbatim}

\hfill\break
\hfill\break
\hfill\break
\hfill\break

\textbf{Question 7. (4 points) Is sampling equal across islands? Create
a plot to visualize the number of anole species sampled from each
island. Be sure to label your axes and add a title.}\\

\begin{Shaded}
\begin{Highlighting}[]
\NormalTok{anolis }\SpecialCharTok{\%\textgreater{}\%} 
  \FunctionTok{ggplot}\NormalTok{(}\FunctionTok{aes}\NormalTok{(}\AttributeTok{x=}\NormalTok{island))}\SpecialCharTok{+}
  \FunctionTok{geom\_bar}\NormalTok{(}\FunctionTok{aes}\NormalTok{(}\AttributeTok{color=}\NormalTok{island))}\SpecialCharTok{+}
  \FunctionTok{labs}\NormalTok{(}\AttributeTok{title =} \StringTok{"Species per Island"}\NormalTok{, }\AttributeTok{x=}\StringTok{"Island"}\NormalTok{,}\AttributeTok{y=}\StringTok{"Number of Species"}\NormalTok{)}
\end{Highlighting}
\end{Shaded}

\pandocbounded{\includegraphics[keepaspectratio]{midterm1_files/figure-latex/unnamed-chunk-7-1.pdf}}
\#Sampling is not equal across islands, Cuba and Hispaniola were sampled
much more than Jamaica or Puerto Rico\\
\strut \\
\strut \\

\textbf{Question 8. (2 points) Which habitat types are represented in
the data? Display the names of the habitat types.}\\
\strut \\

\begin{Shaded}
\begin{Highlighting}[]
\NormalTok{anolis }\SpecialCharTok{\%\textgreater{}\%} 
  \FunctionTok{distinct}\NormalTok{(habitat)}
\end{Highlighting}
\end{Shaded}

\begin{verbatim}
##        habitat
## 1 Trunk-ground
## 2         Twig
## 3   Grass-bush
## 4  Trunk-crown
\end{verbatim}

\hfill\break
\hfill\break
\hfill\break

\textbf{Question 9. (4 points) Is sampling equal across habitat types?
Create a plot to visualize the number of anole species sampled from each
habitat type. Be sure to label your axes and add a title.}\\
\strut \\
\strut \\

\begin{Shaded}
\begin{Highlighting}[]
\NormalTok{anolis }\SpecialCharTok{\%\textgreater{}\%} 
  \FunctionTok{ggplot}\NormalTok{(}\FunctionTok{aes}\NormalTok{(}\AttributeTok{x=}\NormalTok{habitat))}\SpecialCharTok{+}
  \FunctionTok{geom\_bar}\NormalTok{(}\FunctionTok{aes}\NormalTok{(}\AttributeTok{color=}\NormalTok{habitat))}\SpecialCharTok{+}
  \FunctionTok{labs}\NormalTok{(}\AttributeTok{title =} \StringTok{"Species per Habitat"}\NormalTok{, }\AttributeTok{x=}\StringTok{"Habitat"}\NormalTok{,}\AttributeTok{y=}\StringTok{"Number of Species"}\NormalTok{)}
\end{Highlighting}
\end{Shaded}

\pandocbounded{\includegraphics[keepaspectratio]{midterm1_files/figure-latex/unnamed-chunk-9-1.pdf}}
\#Trunk-ground habitat was sampled much more than the others\\
\strut \\

\textbf{Question 10. (4 points) The morphology of anoles varies based on
their habitat. How does the range of hindlimb length compare among
different habitats? Create a plot to visualize the distribution of
hindlimb lengths across habitat types. Be sure to label your axes and
add a title.}\\
\strut \\
\strut \\

\begin{Shaded}
\begin{Highlighting}[]
\NormalTok{anolis }\SpecialCharTok{\%\textgreater{}\%} 
  \FunctionTok{ggplot}\NormalTok{(}\FunctionTok{aes}\NormalTok{(}\AttributeTok{x=}\NormalTok{habitat,}\AttributeTok{y=}\NormalTok{hindlimb\_length\_mm))}\SpecialCharTok{+}
  \FunctionTok{geom\_boxplot}\NormalTok{(}\FunctionTok{aes}\NormalTok{(}\AttributeTok{color=}\NormalTok{habitat))}\SpecialCharTok{+}
  \FunctionTok{labs}\NormalTok{(}\AttributeTok{title =} \StringTok{"Hindlimb Length vs. Habitat"}\NormalTok{, }\AttributeTok{x=}\StringTok{"Habitat"}\NormalTok{,}\AttributeTok{y=}\StringTok{"Hind Limb Length"}\NormalTok{)}
\end{Highlighting}
\end{Shaded}

\pandocbounded{\includegraphics[keepaspectratio]{midterm1_files/figure-latex/unnamed-chunk-10-1.pdf}}
\#Twig anoles have the shortes hind limb length\\
\strut \\

\textbf{Question 11. (4 points) The plot above is compelling, but don't
we expect larger lizards to have longer limbs? What about tail length?
Shouldn't longer lizards have longer tails? To correct for this, make
two new columns: 1. \texttt{ratio\_of\_hindlimb\_to\_body}, and 2.
\texttt{ratio\_of\_tail\_to\_body}. Don't forget to add these columns to
the anolis data frame.}\\
\strut \\

\begin{Shaded}
\begin{Highlighting}[]
\NormalTok{anolis}\SpecialCharTok{$}\NormalTok{ratio\_of\_hindlimb\_to\_body }\OtherTok{\textless{}{-}}\NormalTok{ (anolis}\SpecialCharTok{$}\NormalTok{hindlimb\_length\_mm}\SpecialCharTok{/}\NormalTok{anolis}\SpecialCharTok{$}\NormalTok{body\_length\_mm)}
\NormalTok{anolis}\SpecialCharTok{$}\NormalTok{ratio\_of\_tail\_to\_body }\OtherTok{\textless{}{-}}\NormalTok{ (anolis}\SpecialCharTok{$}\NormalTok{tail\_length\_mm}\SpecialCharTok{/}\NormalTok{anolis}\SpecialCharTok{$}\NormalTok{body\_length\_mm)}
\end{Highlighting}
\end{Shaded}

\hfill\break
\hfill\break
\hfill\break

\textbf{Question 12. (4 points) Create a new plot that examines the
distribution of \texttt{ratio\_of\_hindlimb\_to\_body} across habitat
types. How does this plot differ from the one you made in Problem 10? Be
sure to label your axes and add a title.}\\
\strut \\

\begin{Shaded}
\begin{Highlighting}[]
\NormalTok{anolis }\SpecialCharTok{\%\textgreater{}\%} 
  \FunctionTok{ggplot}\NormalTok{(}\FunctionTok{aes}\NormalTok{(}\AttributeTok{x=}\NormalTok{habitat,}\AttributeTok{y=}\NormalTok{ratio\_of\_hindlimb\_to\_body))}\SpecialCharTok{+}
  \FunctionTok{geom\_boxplot}\NormalTok{(}\FunctionTok{aes}\NormalTok{(}\AttributeTok{color=}\NormalTok{habitat))}\SpecialCharTok{+}
  \FunctionTok{labs}\NormalTok{(}\AttributeTok{title =} \StringTok{"Hindlimb to Bodylength Ratio vs. Habitat"}\NormalTok{, }\AttributeTok{x=}\StringTok{"Habitat"}\NormalTok{,}\AttributeTok{y=}\StringTok{"Hind Limb to Body Length (mm)"}\NormalTok{)}
\end{Highlighting}
\end{Shaded}

\pandocbounded{\includegraphics[keepaspectratio]{midterm1_files/figure-latex/unnamed-chunk-12-1.pdf}}
\#This plot differs in that it shows the proportional hindlimb length
opposed to raw hindlimb length\\
\strut \\
\strut \\

\textbf{Problem 13. (4 points) A longer tail provides better balance and
agility. Create a plot that examines the relationship between body
length and tail length. Color the points by habitat type and add a line
of best fit. What does this plot suggest about the relationship between
body length and tail length? What do you notice about lizards in the
\texttt{Grass-bush} habitat? Be sure to label your axes and add a
title.}

\#The graph suggests that longer bodys are correlated to longer tails,
and also depicts how tail length is longer in proportion to body length
among grass-bush lizards\\
\strut \\

\begin{Shaded}
\begin{Highlighting}[]
\NormalTok{anolis }\SpecialCharTok{\%\textgreater{}\%} 
  \FunctionTok{ggplot}\NormalTok{(}\FunctionTok{aes}\NormalTok{(}\AttributeTok{x=}\NormalTok{body\_length\_mm,}\AttributeTok{y=}\NormalTok{tail\_length\_mm))}\SpecialCharTok{+}
  \FunctionTok{geom\_point}\NormalTok{(}\FunctionTok{aes}\NormalTok{(}\AttributeTok{color=}\NormalTok{habitat))}\SpecialCharTok{+}
  \FunctionTok{geom\_smooth}\NormalTok{(}\AttributeTok{method=}\NormalTok{lm, }\AttributeTok{se=}\NormalTok{T)}\SpecialCharTok{+}
  \FunctionTok{labs}\NormalTok{(}\AttributeTok{title=}\StringTok{"Body vs. Tail Length"}\NormalTok{, }\AttributeTok{x=}\StringTok{"Body Length (mm)"}\NormalTok{, }\AttributeTok{y=}\StringTok{"Tail Length (mm)"}\NormalTok{)}
\end{Highlighting}
\end{Shaded}

\begin{verbatim}
## `geom_smooth()` using formula = 'y ~ x'
\end{verbatim}

\pandocbounded{\includegraphics[keepaspectratio]{midterm1_files/figure-latex/unnamed-chunk-13-1.pdf}}

\hfill\break
\hfill\break

\textbf{Problem 14. (4 points) Toepad lamellae are transverse,
plate-like structures found on the ventral surface of the digits. They
are a key adaptation that allows anoles to cling to and move efficiently
on smooth and vertical surfaces. What is the mean number of toepad
lamellae for each habitat type?}\\

\begin{Shaded}
\begin{Highlighting}[]
\NormalTok{anolis}\SpecialCharTok{\%\textgreater{}\%} 
  \FunctionTok{select}\NormalTok{(habitat,toepad\_lamellae\_count) }\SpecialCharTok{\%\textgreater{}\%} 
  \FunctionTok{filter}\NormalTok{(habitat}\SpecialCharTok{==}\StringTok{"Trunk{-}ground"}\NormalTok{) }\SpecialCharTok{\%\textgreater{}\%} 
  \FunctionTok{summarize}\NormalTok{(}\FunctionTok{mean}\NormalTok{(toepad\_lamellae\_count)) }
\end{Highlighting}
\end{Shaded}

\begin{verbatim}
##   mean(toepad_lamellae_count)
## 1                    29.95652
\end{verbatim}

\begin{Shaded}
\begin{Highlighting}[]
\NormalTok{anolis}\SpecialCharTok{\%\textgreater{}\%} 
  
  \FunctionTok{select}\NormalTok{(habitat,toepad\_lamellae\_count) }\SpecialCharTok{\%\textgreater{}\%} 
  \FunctionTok{filter}\NormalTok{(habitat}\SpecialCharTok{==}\StringTok{"Twig"}\NormalTok{)}\SpecialCharTok{\%\textgreater{}\%} 
  \FunctionTok{summarize}\NormalTok{(}\FunctionTok{mean}\NormalTok{(toepad\_lamellae\_count))}
\end{Highlighting}
\end{Shaded}

\begin{verbatim}
##   mean(toepad_lamellae_count)
## 1                        27.6
\end{verbatim}

\begin{Shaded}
\begin{Highlighting}[]
\NormalTok{anolis}\SpecialCharTok{\%\textgreater{}\%} 
  
  
  \FunctionTok{select}\NormalTok{(habitat,toepad\_lamellae\_count) }\SpecialCharTok{\%\textgreater{}\%} 
  \FunctionTok{filter}\NormalTok{(habitat}\SpecialCharTok{==}\StringTok{"Grass{-}bush"}\NormalTok{)}\SpecialCharTok{\%\textgreater{}\%} 
  \FunctionTok{summarize}\NormalTok{(}\FunctionTok{mean}\NormalTok{(toepad\_lamellae\_count))}
\end{Highlighting}
\end{Shaded}

\begin{verbatim}
##   mean(toepad_lamellae_count)
## 1                     28.3125
\end{verbatim}

\begin{Shaded}
\begin{Highlighting}[]
\NormalTok{anolis}\SpecialCharTok{\%\textgreater{}\%} 
  
  \FunctionTok{select}\NormalTok{(habitat,toepad\_lamellae\_count) }\SpecialCharTok{\%\textgreater{}\%} 
  \FunctionTok{filter}\NormalTok{(habitat}\SpecialCharTok{==}\StringTok{"Trunk{-}crown"}\NormalTok{)}\SpecialCharTok{\%\textgreater{}\%} 
  \FunctionTok{summarize}\NormalTok{(}\FunctionTok{mean}\NormalTok{(toepad\_lamellae\_count))}
\end{Highlighting}
\end{Shaded}

\begin{verbatim}
##   mean(toepad_lamellae_count)
## 1                        38.5
\end{verbatim}

\hfill\break
\hfill\break
\hfill\break
\hfill\break

\textbf{Problem 15. (5 points) The number of toepad lamellae is
significantly different for trunk-crown species. But, is this consistent
across all islands? Make a plot that shows the range in number of toepad
lamellae by island for trunk-crown species only. Be sure to label your
axes and add a title.}\\
\strut \\
\strut \\

\begin{Shaded}
\begin{Highlighting}[]
\NormalTok{anolis }\SpecialCharTok{\%\textgreater{}\%} 
   \FunctionTok{filter}\NormalTok{(habitat}\SpecialCharTok{==}\StringTok{"Trunk{-}crown"}\NormalTok{)}\SpecialCharTok{\%\textgreater{}\%} 
   \FunctionTok{ggplot}\NormalTok{(}\FunctionTok{aes}\NormalTok{(}\AttributeTok{x=}\NormalTok{island,}\AttributeTok{y=}\NormalTok{toepad\_lamellae\_count))}\SpecialCharTok{+}
  \FunctionTok{geom\_boxplot}\NormalTok{(}\FunctionTok{aes}\NormalTok{(}\AttributeTok{color=}\NormalTok{island))}\SpecialCharTok{+}
  \FunctionTok{labs}\NormalTok{(}\AttributeTok{title=}\StringTok{"Trunk{-}Crown Anole Toepad Lamellae \# by Island"}\NormalTok{, }\AttributeTok{x=}\StringTok{"Island"}\NormalTok{, }\AttributeTok{y=}\StringTok{"Toepad Lamellae Count"}\NormalTok{)}
\end{Highlighting}
\end{Shaded}

\pandocbounded{\includegraphics[keepaspectratio]{midterm1_files/figure-latex/unnamed-chunk-15-1.pdf}}

\hfill\break
\hfill\break

\subsection{Submit the Midterm}\label{submit-the-midterm}

\begin{enumerate}
\def\labelenumi{\arabic{enumi}.}
\tightlist
\item
  Save your work and knit the .rmd file.\\
\item
  Open the .html file and ``print'' it to a .pdf file in Google Chrome
  (not Safari).\\
\item
  Go to the class Canvas page and open Gradescope.\\
\item
  Submit your .pdf file to the midterm assignment- be sure to assign the
  pages to the correct questions.\\
\item
  Commit and push your work to your repository.
\end{enumerate}

\end{document}
